\documentclass[10pt]{beamer}
\usetheme{metropolis}
\usepackage{graphicx}
\usepackage{listings}
\usepackage{booktabs}
\usepackage{xcolor}
\usepackage{hyperref}

\definecolor{codebg}{HTML}{F5F5F5}
\definecolor{codegreen}{HTML}{2E7D32}
\definecolor{codegray}{HTML}{757575}
\definecolor{codeblue}{HTML}{1565C0}

\lstdefinestyle{rstyle}{
  language=R, backgroundcolor=\color{codebg}, basicstyle=\ttfamily\scriptsize,
  keywordstyle=\color{codeblue}\bfseries, stringstyle=\color{codegreen},
  commentstyle=\color{codegray}\itshape, breaklines=true, frame=single,
  rulecolor=\color{codegray!30}, showstringspaces=false, tabsize=2,
  morekeywords={library,ggplot,aes,geom_point,geom_col,geom_smooth,
    theme,theme_minimal,theme_classic,theme_bw,theme_void,theme_gray,
    theme_economist,theme_wsj,theme_fivethirtyeight,theme_tufte,
    element_text,element_blank,element_line,element_rect,
    labs,ggsave,margin,unit,rel}
}
\lstdefinestyle{pythonstyle}{
  language=Python, backgroundcolor=\color{codebg}, basicstyle=\ttfamily\scriptsize,
  keywordstyle=\color{codeblue}\bfseries, stringstyle=\color{codegreen},
  commentstyle=\color{codegray}\itshape, breaklines=true, frame=single,
  rulecolor=\color{codegray!30}, showstringspaces=false, tabsize=4,
  morekeywords={import,from,as,plt,rcParams,update,style,use,savefig,
    set_visible,tick_params,set_facecolor,set_xlabel,set_ylabel,set_title,
    tight_layout,subplots,figure,text}
}

\title{Themes \& Publication-Quality Styling}
\subtitle{Workshop 2 --- Module 6}
\author{Prof.Asoc.\ Endri Raco}
\institute{Data Visualization for Data Scientists \\ UNYT Tirana}
\date{2025/26}

\begin{document}

\begin{frame}
  \titlepage
\end{frame}

\begin{frame}{Module Roadmap}
  \begin{enumerate}
    \item What a theme controls: the non-data layer
    \item Five built-in ggplot2 themes
    \item The four \texttt{element\_*()} building blocks
    \item Writing a reusable custom theme function
    \item Publication themes via \texttt{ggthemes}
    \item Python equivalents: \texttt{rcParams}, \texttt{plt.style.use()}
    \item Export formats and resolution for publication
  \end{enumerate}
  \begin{exampleblock}{Learning Outcome}
    You will transform any chart from default styling to publication
    quality using \texttt{theme()} elements and create a reusable
    branded theme function.
  \end{exampleblock}
\end{frame}

\section{What a Theme Controls}

\begin{frame}{theme() Controls Every Non-Data Element}
  \begin{center}
    \includegraphics[width=0.85\textwidth]{figures_w02m06/theme_anatomy}
  \end{center}
  \begin{alertblock}{Pro-Tip}
    A theme modifies \emph{appearance}, never \emph{data}. You can
    change the theme without changing a single data point, aesthetic
    mapping, or statistic. This is the grammar's separation of
    content from presentation.
  \end{alertblock}
\end{frame}

\begin{frame}{Before \& After: The Power of theme()}
  \begin{center}
    \includegraphics[width=0.92\textwidth]{figures_w02m06/before_after_theme}
  \end{center}
\end{frame}

\section{Built-In Themes}

\begin{frame}{Five Core ggplot2 Themes}
  \begin{center}
    \includegraphics[width=0.95\textwidth]{figures_w02m06/builtin_themes}
  \end{center}
  \begin{exampleblock}{Data Insight}
    \texttt{theme\_minimal()} is the recommended starting point: white
    background, subtle gridlines, no box frame. Add \texttt{base\_size}
    for global font scaling: \texttt{theme\_minimal(base\_size = 14)}
    for presentations; \texttt{base\_size = 9} for journal figures.
  \end{exampleblock}
\end{frame}

\begin{frame}[fragile]{Choosing a Base Theme}
  \centering
  \begin{tabular}{lp{5.5cm}}
    \toprule
    \textbf{Theme} & \textbf{Best For} \\
    \midrule
    \texttt{theme\_minimal()} & Default recommendation: clean, modern \\
    \texttt{theme\_classic()} & Journal articles with axis-line style \\
    \texttt{theme\_bw()} & Print-safe with panel border \\
    \texttt{theme\_void()} & Maps, network diagrams, pie charts \\
    \texttt{theme\_gray()} & ggplot2 default (avoid for publication) \\
    \bottomrule
  \end{tabular}

  \vspace{6pt}
  {\small Start with one of these, then customize individual
  elements with \texttt{theme(...)}. Never build a theme from scratch.}
\end{frame}

\section{The element\_*() System}

\begin{frame}{Four Building Blocks}
  \begin{center}
    \includegraphics[width=0.88\textwidth]{figures_w02m06/element_types}
  \end{center}
\end{frame}

\begin{frame}[fragile]{Common theme() Customisations}
\begin{lstlisting}[style=rstyle]
theme_minimal(base_size = 11) +
theme(
  # Title: bold, 14pt, left-aligned
  plot.title = element_text(face = "bold", size = 14, hjust = 0),
  plot.subtitle = element_text(color = "grey50", size = 10),
  plot.caption = element_text(face = "italic", size = 8, hjust = 0),

  # Axes: remove top/right gridlines
  panel.grid.major.x = element_blank(),
  panel.grid.minor = element_blank(),

  # Axis labels: slight margin
  axis.title = element_text(face = "bold", size = 10),
  axis.text = element_text(size = 8),

  # Legend: bottom, horizontal
  legend.position = "bottom",
  legend.title = element_text(face = "bold", size = 9),

  # Margins: breathing room
  plot.margin = margin(t = 10, r = 15, b = 10, l = 10)
)
\end{lstlisting}
\end{frame}

\begin{frame}[fragile]{Removing Elements with element\_blank()}
\begin{lstlisting}[style=rstyle]
# Remove specific elements entirely
theme(
  # Remove y-axis gridlines (for horizontal bar charts)
  panel.grid.major.y = element_blank(),

  # Remove axis ticks and text on one side
  axis.ticks.x = element_blank(),
  axis.text.x = element_blank(),

  # Remove legend
  legend.position = "none",

  # Remove panel border
  panel.border = element_blank(),

  # Remove strip background (facet headers)
  strip.background = element_blank()
)
\end{lstlisting}
  \begin{alertblock}{Pro-Tip}
    \texttt{element\_blank()} is the Tufte eraser: every non-data
    element you remove increases the data-ink ratio (W01-M02).
  \end{alertblock}
\end{frame}

\section{Custom Reusable Themes}

\begin{frame}[fragile]{Writing theme\_unyt(): Your Branded Theme}
  \begin{columns}[T]
    \begin{column}{0.52\textwidth}
\begin{lstlisting}[style=rstyle]
# Define once, use everywhere
theme_unyt <- function(
    base_size = 11) {
  theme_minimal(
    base_size = base_size) +
  theme(
    plot.title = element_text(
      face = "bold", size = rel(1.3)),
    plot.subtitle = element_text(
      color = "grey50"),
    plot.caption = element_text(
      face = "italic", hjust = 0),
    panel.grid.major.x =
      element_blank(),
    panel.grid.minor =
      element_blank(),
    legend.position = "bottom",
    plot.margin = margin(10,15,10,10)
  )
}

# Usage: one line replaces 15+
ggplot(df, aes(x, y)) +
  geom_point() +
  theme_unyt()
\end{lstlisting}
    \end{column}
    \begin{column}{0.45\textwidth}
      \includegraphics[width=\textwidth]{figures_w02m06/custom_theme}

      \vspace{4pt}
      {\small A custom theme function
      ensures brand consistency across
      all charts in a project. Define
      it once in a utility script,
      source it in every analysis.}
    \end{column}
  \end{columns}
\end{frame}

\section{ggthemes: Publication Styles}

\begin{frame}{Mimicking Publication Styles}
  \begin{center}
    \includegraphics[width=0.95\textwidth]{figures_w02m06/ggthemes_gallery}
  \end{center}
  \begin{exampleblock}{Data Insight}
    The \texttt{ggthemes} package provides themes that mimic
    \emph{The Economist}, \emph{WSJ}, \emph{FiveThirtyEight}, and
    Tufte's style. Install with \texttt{install.packages("ggthemes")}.
    These are useful for matching a target publication's visual identity.
  \end{exampleblock}
\end{frame}

\section{Python Equivalents}

\begin{frame}{rcParams: Python's Global Theme System}
  \begin{center}
    \includegraphics[width=0.85\textwidth]{figures_w02m06/rcparams}
  \end{center}
\end{frame}

\begin{frame}[fragile]{Three Python Approaches to Theming}
  \begin{columns}[T]
    \begin{column}{0.48\textwidth}
      \textbf{1. rcParams (global)}
\begin{lstlisting}[style=pythonstyle]
plt.rcParams.update({
    "figure.figsize": (7, 5),
    "figure.dpi": 150,
    "font.size": 11,
    "axes.spines.top": False,
    "axes.spines.right": False,
    "axes.grid": True,
    "grid.alpha": 0.15,
})
\end{lstlisting}

      \vspace{4pt}
      \textbf{2. plt.style.use()}
\begin{lstlisting}[style=pythonstyle]
# Built-in styles
plt.style.use("seaborn-v0_8-whitegrid")
plt.style.use("ggplot")
plt.style.use("bmh")
\end{lstlisting}
    \end{column}
    \begin{column}{0.48\textwidth}
      \textbf{3. Per-chart manual}
\begin{lstlisting}[style=pythonstyle]
fig, ax = plt.subplots()
ax.scatter(x, y)

# Manual Tufte cleanup
ax.spines["top"].set_visible(False)
ax.spines["right"].set_visible(False)
ax.set_facecolor("white")
ax.grid(alpha=0.12)
ax.tick_params(labelsize=8)
ax.set_title("Title",
    fontweight="bold", fontsize=12)
fig.text(0.12, 0.01,
    "Source: ...",
    fontsize=6, color="#AAA")
\end{lstlisting}

      \vspace{4pt}
      {\small \textbf{plotnine} uses
      \texttt{+ theme\_minimal()}
      identical to ggplot2.}
    \end{column}
  \end{columns}
\end{frame}

\section{Export for Publication}

\begin{frame}{Export Format Guide}
  \begin{center}
    \includegraphics[width=0.92\textwidth]{figures_w02m06/export_formats}
  \end{center}
\end{frame}

\begin{frame}[fragile]{Export in R vs Python}
  \begin{columns}[T]
    \begin{column}{0.48\textwidth}
      \textbf{R (ggsave)}
\begin{lstlisting}[style=rstyle]
# Vector (always preferred)
ggsave("chart.pdf",
  plot = p,
  width = 7, height = 5,
  device = cairo_pdf)

# High-res raster
ggsave("chart.png",
  plot = p,
  width = 7, height = 5,
  dpi = 300)

# SVG for web
ggsave("chart.svg",
  plot = p,
  width = 7, height = 5)
\end{lstlisting}
    \end{column}
    \begin{column}{0.48\textwidth}
      \textbf{Python (savefig)}
\begin{lstlisting}[style=pythonstyle]
# Vector (always preferred)
fig.savefig("chart.pdf",
    bbox_inches="tight")

# High-res raster
fig.savefig("chart.png",
    dpi=300,
    bbox_inches="tight")

# SVG for web
fig.savefig("chart.svg",
    bbox_inches="tight")

# ALWAYS close after saving
plt.close()
\end{lstlisting}
    \end{column}
  \end{columns}

  \vspace{4pt}
  \begin{alertblock}{Pro-Tip}
    Always use \texttt{bbox\_inches="tight"} in Python to prevent
    label clipping. In R, \texttt{cairo\_pdf} produces better
    anti-aliased text than the default PDF device.
  \end{alertblock}
\end{frame}

\section{Synthesis}

\begin{frame}{Theme Workflow}
  \begin{enumerate}
    \item Start with a \textbf{base theme}: \texttt{theme\_minimal(base\_size)}
    \item Remove non-data ink with \textbf{element\_blank()}: minor gridlines,
          unnecessary axis ticks, panel borders
    \item Customise text with \textbf{element\_text()}: bold titles, italic
          captions, sized axis labels
    \item Wrap into a \textbf{custom function}: \texttt{theme\_unyt()} for
          brand consistency
    \item Export as \textbf{PDF} (vector) for print, \textbf{PNG 300dpi}
          for screen, \textbf{SVG} for web
    \item Never export as \textbf{JPEG} (lossy compression destroys
          sharp lines and text)
  \end{enumerate}
\end{frame}

\begin{frame}{Key Takeaways}
  \begin{enumerate}
    \item \textbf{theme()} controls all non-data elements: titles, axes,
          gridlines, backgrounds, legends, margins.
    \item Four building blocks: \textbf{element\_text()}, \textbf{element\_line()},
          \textbf{element\_rect()}, \textbf{element\_blank()}.
    \item Start with \textbf{theme\_minimal()} and customise; never build
          from scratch.
    \item Create a \textbf{reusable theme function} for project-wide
          brand consistency.
    \item In Python: \textbf{rcParams} for global settings, manual spine
          removal per chart, or \texttt{plt.style.use()}.
    \item Export: \textbf{PDF/SVG} (vector) > \textbf{PNG 300dpi} (raster)
          > never JPEG.
  \end{enumerate}
\end{frame}

\begin{frame}{Essential Readings}
  \begin{itemize}
    \item Wickham, H. (2016). \emph{ggplot2}, Chapter 8 (Themes).
    \item Wilke, C.~O. (2019). \emph{Fundamentals of Data Visualization},
          Chapter 22 (Titles, Captions, and Tables).
    \item Chang, W. (2024). \emph{R Graphics Cookbook}, 2nd ed., Chapter 9
          (Controlling the Overall Appearance).
    \item ggthemes documentation: \texttt{https://jrnold.github.io/ggthemes/}
  \end{itemize}
  \vspace{6pt}
  \textbf{Next Module}: seaborn \& plotnine for Python (W02-M07)
\end{frame}

\end{document}
